\itemhead{1.2 例题 · 例题6（教材p.20）}{XJ-B1-C01-EX-P020-Y0562}
\label{XJ-B1-C01-EX-P020-Y0562}
{\footnotesize 湘教版教材p.20 / PDF第27页 / 1.2 例题 / 例题6}
\field{原题}{指出下列命题中使用了什么量词以及量词的作用范围，并把量词用相应的数学符号取代：
\subq{1}对任意正实数 $a$，$a^2-a-2>0$；
\subq{2}对某个大于 10 的正整数 $n$，$(\sqrt2)^n=1024$。}
\field{参考答案}{（1）全称量词，范围为 $\mathbb R_+$，写作 $\forall a\in\mathbb R_+,a^2-a-2>0$；（2）存在量词，范围为大于10的正整数，写作 $\exists n\in\mathbb Z_+,n>10,(\sqrt2)^n=1024$。}
{\small \field{关联知识}{L03 全称、存在量词及真假的判定（教材帮p.38、39 / PDF47、48）}
\field{方法与易错点}{先确认量词、作用对象及论域，再采用任意性证明或实例反例。 易错：用几个满足的数证明全称命题；把“存在”理解为“唯一存在”。}
\field{答案依据}{例题原解位置：湘教版教材p.20及题库原图；独立编写可编辑参考解答。答案册此部分主要收练习、习题，不伪造例题答案册位置。}}
